\documentclass[12pt, a4paper]{article}

% Required Packages
\usepackage{amsmath}
\usepackage{amssymb}
\usepackage{geometry}
\usepackage{graphicx}
\usepackage{hyperref}

% Page Geometry Setup
\geometry{
    top=2.5cm,
    bottom=2.5cm,
    left=2.5cm,
    right=2.5cm
}

\title{\textbf{Mathematical Formulation of the Microgrid Dispatch Problem with Non-Linear Battery Degradation}}
\author{Optimization Project Formulation}
\date{\today}

\begin{document}

\maketitle

\section{Problem Definition}
We are solving a Finite-Horizon Discrete-Time Optimal Control Problem. The goal is to determine the optimal power dispatch of a stationary battery storage system over a finite horizon to minimize the total operational cost. The cost comprises the net grid energy exchange, grid usage fees, and a non-linear penalty representing battery cycle degradation.

\section{Sets and Indices}
\begin{itemize}
    \item $\mathcal{T} = \{1, 2, \dots, N\}$: The optimization horizon.
    \item $t \in \mathcal{T}$: A discrete time step index, where the duration of each step is $\Delta t$ (e.g., $\Delta t = 1$ hour).
    \item $\mathcal{S} = \{1, 2, \dots, S\}$: The set of price scenarios, representing intraday price uncertainty relative to the day-ahead market.
    \item $s \in \mathcal{S}$: Scenario index.
\end{itemize}

\section{Input Data (Parameters)}

\subsection{System Parameters}
\begin{itemize}
    \item $E^{max}$: Maximum energy capacity of the battery (kWh).
    \item $E^{min}$: Minimum allowable energy capacity to prevent deep depletion (kWh).
    \item $P^{max}$: Maximum allowable charge and discharge power (kW).
    \item $\eta_{c}$: Charging efficiency ($0 < \eta_{c} \le 1$).
    \item $\eta_{d}$: Discharging efficiency ($0 < \eta_{d} \le 1$).
    \item $\alpha$: Degradation penalty coefficient (\$/kW$^2$). Translates physical stress into economic cost.
    \item $\lambda^{grid}$: Grid usage fee for network access and transaction costs (\$/kWh).
    \item $E_0$: Initial battery state of energy at $t=0$ (kWh).
\end{itemize}

\subsection{Time-Series Parameters}
\begin{itemize}
    \item $P_{load,t}$: Household power demand at time $t$ (kW). Note: This is currently a point estimate from the ML prediction service.
    \item $P_{pv,t}$: Solar photovoltaic generation at time $t$ (kW).
    \item $\lambda^{DA}_t$: Day-Ahead electricity price at time $t$ (\$/kWh), fetched from ENTSO-E.
    \item $\epsilon_{t,s}$: Intraday price spread (stochastic noise) at time $t$ under scenario $s$ (\$/kWh), simulated via an AR(1)-GARCH(1,1) process to capture volatility clustering.
    \item $\lambda^{buy}_{t,s}$: Price of purchasing electricity from the grid at time $t$ under scenario $s$ (\$/kWh). Defined as $\lambda^{buy}_{t,s} = \lambda^{DA}_t + \epsilon_{t,s}$.
    \item $\lambda^{sell}_{t,s}$: Price of selling electricity to the grid at time $t$ under scenario $s$ (\$/kWh). Defined as $\lambda^{sell}_{t,s} = 0.9 \cdot \lambda^{DA}_t + \epsilon_{t,s}$, which incorporates the 10\% transaction/selling discount.
\end{itemize}

\section{Decision Variables}
These are the continuous variables optimized by the Interior Point solver at each time step $t \in \mathcal{T}$:
\begin{itemize}
    \item $P_{buy,t} \in [0, 20]$: Power imported from the grid (kW).
    \item $P_{sell,t} \in [0, 20]$: Power exported to the grid (kW).
    \item $P_{ch,t} \in \mathbb{R}_{\ge 0}$: Power sent to the battery for charging (kW).
    \item $P_{dis,t} \in \mathbb{R}_{\ge 0}$: Power drawn from the battery for discharging (kW).
    \item $E_t \in \mathbb{R}_{\ge 0}$: Battery State of Energy at the end of time step $t$ (kWh).
\end{itemize}

\section{Mathematical Formulation}

\subsection{Objective Function}
The objective is to minimize the \textbf{expected total cost} $\mathbb{E}[J]$ over all price scenarios:

\begin{equation}
    \min \mathbb{E}[J] = \frac{1}{S} \sum_{s \in \mathcal{S}} \sum_{t=1}^{N} \left[ \left( \lambda^{buy}_{t,s} P_{buy,t} - \lambda^{sell}_{t,s} P_{sell,t} + \lambda^{grid} (P_{buy,t} + P_{sell,t}) \right) \Delta t + \alpha \left( P_{ch,t}^2 + P_{dis,t}^2 \right) \right]
\end{equation}

\vspace{0.5cm}
\noindent \textbf{Note on Grid Usage Fees:} The linear term $\lambda^{grid}(P_{buy,t} + P_{sell,t})$ represents real-world costs associated with using the electricity grid, such as network access fees or transaction spreads. Mathematically, this term also acts as a penalty for simultaneous buying and selling: even if $\lambda^{buy} = \lambda^{sell}$, any non-zero flow through the grid incurs a cost, ensuring the solver selects the most efficient net energy exchange.

\noindent \textbf{Note on Non-Linearity:} The quadratic term $\alpha(P_{ch,t}^2 + P_{dis,t}^2)$ models the non-linear degradation proxy. In battery physics, $I^2R$ heating and high C-rates accelerate component wear. Squaring the power mathematically penalizes sudden, aggressive power spikes, forcing the solver to prefer slow, steady charging profiles.

\subsection{Constraints}

\noindent \textbf{1. Power Node Balance:} \\
At every time step, the total power entering the household system must perfectly equal the power leaving or being consumed.
\begin{equation}
    P_{pv,t} + P_{buy,t} + P_{dis,t} = P_{load,t} + P_{sell,t} + P_{ch,t} \quad \forall t \in \mathcal{T}
\end{equation}

\noindent \textbf{2. Battery Energy Dynamics:} \\
The energy state at time $t$ depends on the energy at $t-1$ plus the net energy flowing in/out, adjusted for conversion inefficiencies.
\begin{equation}
    E_t = E_{t-1} + \left( P_{ch,t} \eta_{c} - \frac{P_{dis,t}}{\eta_{d}} \right) \Delta t \quad \forall t \in \mathcal{T}
\end{equation}

\noindent \textbf{3. Operational Power Bounds:} \\
The solver cannot command the battery to charge or discharge faster than its physical limits.
\begin{align}
    0 \le P_{ch,t} &\le P^{max} \quad \forall t \in \mathcal{T} \\
    0 \le P_{dis,t} &\le P^{max} \quad \forall t \in \mathcal{T}
\end{align}

\noindent \textbf{4. State of Energy Bounds:} \\
The battery must not be overcharged or deeply discharged.
\begin{equation}
    E^{min} \le E_t \le E^{max} \quad \forall t \in \mathcal{T}
\end{equation}

\noindent \textbf{5. Grid Interaction Bounds:} \\
Power exchanges with the grid are limited by physical infrastructure (e.g., 20 kW fuse).
\begin{equation}
    0 \le P_{buy,t} \le 20, \quad 0 \le P_{sell,t} \le 20 \quad \forall t \in \mathcal{T}
\end{equation}

\subsection{Consolidated Optimization Problem}
Bringing the objective and all constraints together, the complete stochastic microgrid dispatch problem is formulated as follows:

\begin{equation}
\begin{aligned}
\min_{\mathbf{P}, \mathbf{E}} \quad & \frac{1}{S} \sum_{s \in \mathcal{S}} \sum_{t=1}^{N} \left[ \left( \lambda^{buy}_{t,s} P_{buy,t} - \lambda^{sell}_{t,s} P_{sell,t} + \lambda^{grid}(P_{buy,t} + P_{sell,t}) \right) \Delta t + \alpha \left( P_{ch,t}^2 + P_{dis,t}^2 \right) \right] \\
\text{s.t.} \quad & P_{pv,t} + P_{buy,t} + P_{dis,t} = P_{load,t} + P_{sell,t} + P_{ch,t} && \forall t \in \mathcal{T} \\
& E_t = E_{t-1} + \left( P_{ch,t} \eta_{c} - \frac{P_{dis,t}}{\eta_{d}} \right) \Delta t && \forall t \in \mathcal{T} \\
& 0 \le P_{ch,t} \le P^{max} && \forall t \in \mathcal{T} \\
& 0 \le P_{dis,t} \le P^{max} && \forall t \in \mathcal{T} \\
& E^{min} \le E_t \le E^{max} && \forall t \in \mathcal{T} \\
& 0 \le P_{buy,t} \le 20, \quad 0 \le P_{sell,t} \le 20 && \forall t \in \mathcal{T}
\end{aligned}
\end{equation}

Where $\mathbf{P}$ and $\mathbf{E}$ represent the vectors of all power and energy decision variables over the horizon $\mathcal{T}$.

\section{Algorithm Justification}
Because the objective function contains a convex non-linear term (the squared battery power limits) subject to linear equality and inequality constraints, the problem constitutes a Non-Linear Programming (NLP) model. This requires a specialized Interior Point algorithm, such as IPOPT (Interior Point Optimizer), to reliably compute the global minimum.

\end{document}
